\documentclass[11pt]{article}
\usepackage[T1]{fontenc}
\usepackage{lmodern,amsmath,amssymb,amsthm,microtype}
\usepackage[letterpaper,margin=0.93in]{geometry}
\usepackage[hidelinks]{hyperref}
\hypersetup{pdftitle={The sum of the reciprocal Mersenne cubes is irrational},
  pdfauthor={twinheterodox}}
\newtheorem{theorem}{Theorem}
\newtheorem{lemma}[theorem]{Lemma}
\newtheorem{corollary}[theorem]{Corollary}
\newcommand{\Q}{\mathbb Q}
\newcommand{\Z}{\mathbb Z}
\newcommand{\C}{\mathbb C}
\newcommand{\Ord}{\operatorname{ord}}
\newcommand{\cl}{\mathcal L}
\newcommand{\ph}{\varphi}
\newcommand{\ep}{\varepsilon}
\DeclareMathOperator{\rank}{rank}
\allowdisplaybreaks[1]
\raggedbottom

\title{\Large\bfseries The sum of the reciprocal Mersenne cubes is irrational}
\author{twinheterodox\\[2pt]\normalsize Prepared with Codex}
\date{18 November 2026}
\makeatletter
\renewcommand{\@maketitle}{%
  \begin{center}
    {\@title\par}\medskip
    \begin{tabular}{c}\@author\end{tabular}\par\smallskip
    {\small\@date\par}
  \end{center}}
\makeatother

\begin{document}
\maketitle

\begin{abstract}
The sum of the reciprocal cubes of the Mersenne numbers \(2^n-1\) is
irrational. More generally, for \(|b|>1\) let
\(M_j(b)=\sum_{n\ge1}(b^n-1)^{-j}\).
For every integer \(b\ge2\), we prove that \(1,M_1(b),M_2(b),M_3(b)\)
are linearly independent over \(\Q\). These four numbers are also
linearly independent over \(\Q(b)\) whenever \(b\) is an algebraic
integer with \(|b|>1\) and every conjugate other than \(b\) and
\(\bar b\) lies in the closed unit disk. This includes all Pisot and
Salem numbers.
\end{abstract}

The reciprocals of the Mersenne numbers \(1,3,7,15,\ldots\) sum to
the Erd\H{o}s--Borwein constant. The number \(1\), this constant, and
the sum of the squares of those reciprocals are linearly independent
over \(\Q\) \cite[Theorem~1.3]{PVA}. We prove that adjoining the sum
of their cubes raises the dimension of the rational span from three
to four. For a complex base \(b\) with \(|b|>1\) and a positive
integer \(j\), write
\[
                  M_j(b)=\sum_{m\ge1}\frac{1}{(b^m-1)^j}.
\]

\begin{theorem}\label{thm:main}
For every integer \(b\ge2\), the four numbers
\(1,M_1(b),M_2(b),M_3(b)\) are linearly independent over \(\Q\).
In particular \(1+3^{-3}+7^{-3}+15^{-3}+\cdots\) is irrational.
\end{theorem}

At the golden ratio \(b=(1+\sqrt5)/2\), even coefficients in
\(\Q(\sqrt5)\) give no nontrivial relation among these four numbers.
Corollary~\ref{cor:field} gives the general condition on the
conjugates of an algebraic base.

To compare the theorem with earlier results, consider the Lambert sums
\[
                     Z_s(b)=\sum_{n\ge1}\frac{n^{s-1}}{b^n-1}.
\]
Expanding \(x^j/(1-x)^j\) at \(x=b^{-m}\) and summing in \(m\) gives
\begin{equation}\label{eq:triangle}
 M_1=Z_1,\qquad M_2=Z_2-Z_1,\qquad
 2M_3=Z_3-3Z_2+2Z_1 .
\end{equation}
This triangular change of variables is invertible over \(\Q\).
The first two identities translate the independence of \(1,Z_1,Z_2\),
proved by Postelmans and Van Assche and also by Zudilin
\cite[Theorem~2]{ZudPair}, into the three-number result used above.
The public inventory for the weight-three determinant reports a
companion result asserting independence of \(1,Z_2,Z_3\)
\cite[entry FUT-LI]{UpliftCompanion}.
Both three-number statements are compatible with \(Z_3-3Z_2+2Z_1\) being rational,
which by \eqref{eq:triangle} would make the reciprocal cube sum rational.

We argue by contradiction. For a fixed nonzero triple \(A,B,C\), let
\[
                   \xi=A Z_1(b)+B Z_2(b)+C Z_3(b).
\]
Under a rational relation we can take \(A,B,C,\xi\) all integral.
We adapt the weight-three determinant of \cite[\S\S1--2]{Uplift} to
this combination of Lambert sums. After clearing its denominators and
dividing out a common power of \(b\), we obtain an integer. We show
that its absolute value tends to zero with the size of the determinant
and that infinitely many of these integers are nonzero. The weight-one
term \(AZ_1\) requires a subtraction at zero in each summand. That
subtraction affects both the common power of \(b\) and the estimate
at the value \(\xi\).

\section{One determinant for a quadratic weight}

The series converge absolutely for complex \(|b|>1\). Put \(q=1/b\).
For \(j\ge1\), summing \(x^j\) over the points \(q^m\), \(m\ge1\),
gives \(1/(b^j-1)\). The Euler derivative \(\theta=x\,d/dx\)
multiplies \(x^j\) by \(j\), so \(A+B\theta+C\theta^2\) supplies
the numerator \(A+Bj+Cj^2\).

For a general holomorphic function \(f\), the term \(Af(q^m)\)
tends to \(Af(0)\). To obtain a convergent series we subtract
this limit in every summand and put
\begin{equation}\label{eq:L}
 L(f)=\sum_{m\ge1}
 \bigl[A(f(q^m)-f(0))+B(\theta f)(q^m)+C(\theta^2f)(q^m)\bigr].
\end{equation}
For \(f\) holomorphic on a neighborhood of \(|x|\le |q|\), each
summand is \(O_f(|q|^m)\) for fixed \(A,B,C,q\).
Thus \(L(1)=0\), and \(L(x^j)=(A+Bj+Cj^2)/(b^j-1)\) for \(j\ge1\).
To see why poles at \(b,b^2,\ldots\) are useful, write the first
\(a\) summands at a variable base \(t\) as \(S_a(t)\), where
\[
 h(v)=\frac{A}{v-1}+\frac{Bv}{(v-1)^2}
                +\frac{Cv(v+1)}{(v-1)^3},\qquad
 S_a(t)=\sum_{m=1}^a h(t^m).
\]
For example, \((A,B,C)=(2,-3,1)\) gives \(h(v)=2/(v-1)^3\)
and \(\xi=2M_3(b)\). A general triple lets us test every putative relation.
Expanding the geometric series gives \(\xi=\sum_{m\ge1}h(b^m)\).
In \eqref{eq:L}, replacing \(f(x)=(1-x)^{-1}\) by
\((1-q^a x)^{-1}\) shifts the summation index from \(m\) to \(m+a\).
Only the first \(a\) terms of the series for \(\xi\) are lost. Thus
\[
                   L((1-q^a x)^{-1})=\xi-S_a(b).
\]
The points \(q^m\) where \(L\) samples a function lie inside the unit
disk. The poles \(b^a\) of these test fractions lie outside it, and
each fraction introduces the same unknown \(\xi\).

To clear denominators before evaluating at the actual base and series
value, replace \(b,\xi\) by indeterminates \(t,X\) and treat \(A,B,C\)
as independent indeterminates as well.
For \(K\ge1\) set
\[
 F(x)=\prod_{a=1}^K(x-t^a),\qquad
 D(x)=\prod_{a=1}^K(1-x/t^a).
\]
Define \(\cl_{t,X}\) on fractions with denominator dividing \(D\)
by partial fractions, including the polynomial quotient. Assign
\[
 \mu_0=0,\qquad
 \cl_{t,X}(x^j)=\mu_j=\frac{A+Bj+Cj^2}{t^j-1}\quad(j\ge1),\qquad
 \cl_{t,X}((1-x/t^a)^{-1})=X-S_a(t).
\]
The determinant of these moments is
\begin{equation}\label{eq:delta}
 \Delta_K(t,X)=
 \det\left[\cl_{t,X}\left(\frac{x^{i+j}}{D(x)}\right)\right]_{0\le i,j<K}.
\end{equation}
We suppress \(A,B,C\) in \(\Delta_K(t,X)\). At \(t=b,X=\xi\)
its entries are the actual values of \(L\). For example, take \(K=1\)
and the cube mixture \((A,B,C)=(2,-3,1)\). The determinant has just
one entry, \(X-h(t)\). Multiplication by \((t-1)^3\) gives the
integer polynomial \((t-1)^3 X-2\). At \(t=2,\ X=2M_3(2)\), this
is twice the tail of the cube sum after its first term \(1\).

The geometric poles, the moving denominator \(D\), and the use of
a second Euler derivative follow the pure weight-three construction
\cite[\S\S1--2]{Uplift}. Rational modifications of moment determinants
are treated by Krattenthaler \cite[\S1]{Kratt}, and Zudilin uses
geometric-node Hankel decay for a generalized \(q\)-logarithm
\cite[\S4]{ZudHankel}. The arguments below give the mixed construction
directly, including its weight-one coefficient.

\section{How much of the base divides the determinant}

For large \(K\) the determinant vanishes to high order at \(t=0\).
A lower bound for that order tells us how many powers of the base we
can divide out after clearing the denominators. Write
\[
 g_K(k)=\binom{(K-k)_++1}{2},\qquad
 e_K=\frac{K^2(K+1)}2-\sum_{i=0}^{K-1}g_K(2i),
 \qquad u_+=\max(u,0).
\]

\begin{lemma}\label{lem:base}
As a rational function of \(t\) over \(\Q(A,B,C,X)\),
\(\Ord_0\Delta_K\ge e_K\), where \(e_K\ge0\) and
\(e_K=5K^3/12+O(K^2)\).
\end{lemma}
\begin{proof}
Begin with the monic denominator \(F\). Dividing \(x^k\) by \(F\)
gives a polynomial quotient with coefficients in \(\Z[t]\). Its
moment is regular at zero because every \(\mu_j\) is. So are the
finite tails \(S_a\), including their constant term \(-aA\).
The moment of the proper part is a sum of terms
\[
              -\frac{t^{a(k-1)}}{F'(t^a)}(X-S_a(t)),\qquad 1\le a\le K ,
\]
since \((x-t^a)^{-1}=-t^{-a}(1-x/t^a)^{-1}\).
Summing the orders \(\min(a,c)\) of the root differences
for \(c\ne a\) gives \(\Ord_0 F'(t^a)=Ka-a(a+1)/2\).
The displayed prefactor therefore has order
\(a(a-1)/2+a(k-K)\), bounded below by \(-g_K(k)\).
This includes \(k=0\), where the negative exponent in \(t^{a(k-1)}\)
matters.

Consequently a term indexed by a permutation \(\sigma\) in the
monic-denominator determinant loses at most
\(\sum_i g_K(i+\sigma(i))\) powers. The differences
\(g_K(k+1)-g_K(k)\) are nondecreasing. If \(i<i'\) and \(j<j'\), then
\[
 g_K(i+j)+g_K(i'+j')\ge g_K(i+j')+g_K(i'+j).
\]
Removing inversions gives the largest loss at the identity permutation.
Restoring \(D\) multiplies each entry by \((-1)^K t^{K(K+1)/2}\),
giving the asserted order. Each summand of the loss is at most
\(\binom{K+1}{2}\), so \(e_K\ge0\). Summing squares on an arithmetic
progression gives
\[
  \sum_{0\le i<K/2}\binom{K-2i+1}{2}=K^3/12+O(K^2).\qedhere
\]
\end{proof}

\section{Clearing the cyclotomic denominators}

The remaining denominators vanish at roots of unity. At a primitive
\(d\)-th root, the nodes \(t^a\) with indices congruent modulo \(d\)
coincide. When \(d>K/2\), at most two nodes coincide at any point.
This accounts for the two ranges in the multiplier below.
Let \(\Phi_d\) be the \(d\)-th cyclotomic polynomial and put
\begin{equation}\label{eq:clearer}
 H_K(t)=\prod_{d=1}^K\Phi_d(t)^{c_d},\qquad
 c_d=\begin{cases}
 \lceil K^2/(4d)+2K+17d/16\rceil,&2d\le K,\\
 5(K-d)+3,&2d>K.
 \end{cases}
\end{equation}
The degree \(\sum_{d\le K}c_d\varphi(d)\), where \(\varphi\)
is Euler's totient, controls the size of this product at \(b\).
On the scale \(K^3\log|b|\), division by the base power saves \(5/12\),
and the triple zeros in Section~\ref{sec:decay} will save another
\(1/9\). The multiplier must cost less than their sum. Its exponents
come from the sharpened pure weight-three estimate
\cite[\S4 and Appendix A]{Uplift}.
We show why the same exponents suffice for arbitrary \(A,B,C\).

\begin{lemma}\label{lem:clear}
The polynomial \(P_K=t^{-e_K}H_K\Delta_K\) belongs to
\(\Z[t,A,B,C,X]\). Moreover
\[
 \deg H_K\le\gamma K^3+O(K^2\log(2K)),\qquad
 \gamma=\frac{321}{64\pi^2},\qquad \sum_{d\le K}c_d=O(K^2\log(2K)).
\]
\end{lemma}
\begin{proof}
Fix a primitive \(d\)-th root \(\omega\), put \(\tau=t-\omega\), and
work locally with \(A,B,C,X\) independent. A cluster \(I_r\) consists
of the \(a\) in \(1,\ldots,K\) congruent to \(r\bmod d\).
Write \(m=|I_r|,\ F_r=\prod_{a\in I_r}(x-t^a)\), and use the basis
\((F/F_r)(x-\omega^r)^j,\ 0\le j<m\), over the clusters.
The determinant of this change of basis is nonzero at \(t=\omega\),
so the order we are estimating is preserved. Modulo \(\tau\), each group vanishes
on the other primary factors and gives powers times a unit on its
own factor \((x-\omega^r)^m\).

In a product of two basis elements divided by \(D\), the polynomial
quotient has integral local coefficients and degree at most \(K-2\).
Its moment has at most a simple pole. Between different clusters
the product is divisible by \(F\), so it has no proper part. Within
one cluster the residue numerator at \(t^a\) vanishes to order at least
\(j+k\), and \(F'(t^a)\) has order \(m-1\). Since \(X-S_a\) has
at most a triple pole, the proper-part moment has order at least
\(j+k-m-2\).

For \(2d\le K\), scale each row and column of local index \(j\) by
\(\tau^{w_j}\), where \(w_j=\max(1/2,(m+2)/2-j)\). Half powers mean
that we adjoin a square root of \(\tau\). The sum of two weights pays
for the polynomial part and, within a cluster, for the proper part.
The cost in determinant order for one cluster is
\[
 2\sum_{j=0}^{m-1}w_j=\sum_{j=0}^{m-1}\max(1,m+2-2j)
                  \le m^2/4+2m+1.
\]
The cluster sizes are \(\lfloor K/d\rfloor\) and \(\lceil K/d\rceil\),
with sum \(K\) and sum of squares at most \(K^2/d+d/4\).
This gives the low-order exponent in \eqref{eq:clearer}.

For \(2d>K\), put \(u=K-d\). The \(u\) double clusters are
\(\{a,a+d\}\) with \(1\le a\le u<d\). Use \(V=F/F_r\) and
\(V(x-t^{a+d})\) in each, another local unit change of basis.
Any proper-part entry involving the second vector loses its upper
residue exactly. At the lower node its zero of local order one
cancels the possible simple residue denominator, and \(S_a\) is
regular there since \(a<d\). The entry from two copies of \(V\)
has at most a fourth-order pole. Weights \(2,0\) thus suffice for
the proper parts.
The singleton \(a=d\) needs weight \(3/2\), and the other singletons
all have \(a<d\). The total cost so far is \(4u+3\).

The sole possible pole among polynomial moments here is \(\mu_d\).
Its coefficient form on polynomials \(f,g\) of degree less than \(K\)
is the coefficient of \(x^d\) in the polynomial quotient of \(fg/D\).
If \(\deg f\le d\), the quotient has degree at most \(d-1\) for every
\(g\) in this space. All such \(f\) therefore lie in its radical.
The coefficient form is locally integral and has rank at most
\((K-d-1)_+\) over the fraction field. Row and column scalings preserve
this rank. The scaled full matrix therefore has the form \(U+\tau^{-1}V_1\)
with \(U,V_1\) locally integral and \(\rank V_1\le(K-d-1)_+\).
Multilinearity in the columns costs at most that many extra inverse
powers. Thus \(4u+3+(u-1)_+\le5u+3\), as required.

To return to integer coefficients, put \(R=\Z[A,B,C,X]\).
Monic division by \(F\), followed by partial fractions, introduces
only powers of \(t\) and factors \(t^j-1\) with \(1\le j\le K\)
in the denominators. The displayed moments have the same property.
The local estimates therefore cancel every cyclotomic denominator.
Dividing by the monic cyclotomic factors in \(R[t]\) gives
\(H_K\Delta_K\in R[t,t^{-1}]\). Lemma~\ref{lem:base} removes the
indicated power of \(t\).

For the degree we use the cyclotomic counting method of
Bundschuh and Zudilin \cite[\S3]{BZ}. Summing
\(\ph(d)/d=\sum_{k\mid d}\mu(k)/k\), with \(\mu\) the M\"obius function,
gives \(\sum_{d\le y}\ph(d)/d=6y/\pi^2+O(\log(2y))\).
Partial summation gives the corresponding sums of \(\ph(d)\) and
\(d\ph(d)\) with main terms \(3y^2/\pi^2\) and \(2y^3/\pi^2\)
and errors \(O(y\log(2y))\) and \(O(y^2\log(2y))\).
Substitution in \(\sum c_d\ph(d)\) gives leading contributions
\(161/(64\pi^2)\) and \(5/(2\pi^2)\) times \(K^3\) in the two ranges.
The ceilings contribute \(O(K^2)\). Summing \(c_d=O(K^2/d)\)
gives the last bound.
\end{proof}

\begin{samepage}
The rank correction can already be seen when \(K=5\) and \(d=3\).
The clusters are \(\{1,4\}\), \(\{2,5\}\), and \(\{3\}\), giving
the proper parts a pole-order bound of \(4\cdot2+3=11\).
Among products of two monomials of degree at most four, only
\(x^4x^4\) yields a quotient of degree three after division by \(D\).
The correction therefore allows one more power. This extra power
occurs for \((A,B,C,X)=(3,-2,5,7)\), where the full determinant has
pole order \(12\) at a primitive cube root of unity.\footnote{From the
directory containing the Mersenne verification programs~\cite{MersenneCode},
run \texttt{python3 -B independent-checks.py}.}\par
\end{samepage}

\section{Decay at the value of the series}\label{sec:decay}

We now estimate the determinant at \(t=b,X=\xi\). Consider one
test polynomial \(f(x)=(x-q)^3\). Its first two Euler derivatives
vanish at \(q\), and \(f(q)-f(0)=q^3\), so the \(m=1\) summand
in \eqref{eq:L} is \(Aq^3\). In the determinant we repeat this
construction at \(q,q^2,\ldots\). The same roots that kill the
derivatives also make the surviving value at zero small.

\begin{lemma}\label{lem:decay}
Fix \(|b|>1\) and complex \(A,B,C\). At the true value \(\xi\),
\[
 |P_K(b,A,B,C,\xi)|\le
 \exp\{-\ep K^3\log |b|+O_{b,A,B,C}(K^2\log(2K))\},
 \qquad \ep=\frac{19}{36}-\frac{321}{64\pi^2}>0 .
\]
\end{lemma}
\begin{proof}
Put \(r=|q|\) and, for \(j\ge0\), take
\[
 B_j(x)=\prod_{\ell=1}^j(x-q^{\lceil\ell/3\rceil}),\quad
 s_j=\sum_{\ell=1}^j\lceil\ell/3\rceil,\quad R_j=\lceil j/3\rceil .
\]
These polynomials are monic, so changing both rows and columns
preserves \(\Delta_K\). For one entry write \(R=\max(R_i,R_j)\),
\(W=1/D\), and \(h=B_iB_jW\). On \(|x|\le1\), \(W\) is bounded by
\(\prod_{a\ge1}(1-r^a)^{-1}\). Cauchy's inequalities bound \(W,W',W''\)
on \(|x|\le r\) uniformly in \(K\).

If \(1\le m<R\), at least one of \(B_i,B_j\) contains the factor
\((x-q^m)^3\). The whole summand in \eqref{eq:L} is then
\(-A h(0)\), of modulus \(|A|r^{s_i+s_j}\).
There are fewer than \(K\) such terms. For \(m\ge\max(1,R)\) and on
the entire segment from 0 to \(q^m\), each factor \(x-q^a\) has
modulus at most \(2r^a\). One ordinary derivative of a factor
costs at most \(r^{-R}\) relative to this product bound. Thus
\[
 |h(q^m)-h(0)|\le C_r(i+j+1)\,2^{i+j}r^{s_i+s_j}r^{m-R},
\]
by integrating \(h'\) along the segment. For \(\theta h\) and
\(\theta^2 h\), each term involving a differentiated root factor
gains at least \(r^{m-R}\), since either positive Euler derivative
of \(x-q^a\) equals \(x\). If only \(W\) is differentiated, its
positive Euler derivatives are \(O_r(r^m)\). There are at most
\((i+j+1)^2\) ordered choices of factors for two derivatives.
These estimates include \(R=0\). Summation over \(m\) proves
\[
 |L(h)|\le C_{r,A,B,C}(K+1)^2\,2^{i+j}r^{s_i+s_j}.
\]
Factor \(2^i r^{s_i}\) from row \(i\) and \(2^j r^{s_j}\) from
column \(j\). Leibniz's formula costs a further
\(K![C_{r,A,B,C}(K+1)^2]^K\). Since
\(2\sum_{j<K}s_j=K^3/9+O(K^2)\), this proves
\(\log|\Delta_K(b,\xi)|\le -(K^3/9)\log|b|+O_{b,A,B,C}(K^2)\),
with the usual interpretation at zero.

For fixed \(|b|>1\), the cyclotomic product formula gives
\(\log|\Phi_d(b)|=\ph(d)\log|b|+O_b(1)\) uniformly in \(d\).
Indeed the error is bounded by
\(\sum_{h\ge1}-\log(1-|b|^{-h})<\infty\).
In units of \(K^3\log|b|\), the triple zeros supply \(1/9\) of
decay and division by the common power of \(b\) supplies another
\(5/12\) by Lemma~\ref{lem:base}. The cyclotomic multiplier costs
at most \(\gamma\) by Lemma~\ref{lem:clear}. Their difference
\(1/9+5/12-\gamma=\ep\) is positive already from \(\pi>31/10\).
\end{proof}

\section{Nonvanishing for arbitrarily large sizes}

The coefficients \(A,B,C\) may have either sign in a rational
relation, so we need a nonvanishing argument that also works without
positivity. We encode the moments in a meromorphic function. A
singular determinant of size \(k\) imposes interpolation conditions
at the exterior points \(b,\ldots,b^k\).
If all large sizes were singular, successive approximations would
coincide. The function's poles at the interior points \(q,q^2,\ldots\)
will rule this out. We use the interpolation argument of
\cite[\S5]{Uplift}, in the complex-node form of \cite[\S5]{Golden}.

\begin{lemma}\label{lem:nonzero}
If \(|b|>1\) and \((A,B,C)\ne(0,0,0)\), then
\(\Delta_K(b,\xi)\ne0\) for arbitrarily large \(K\).
\end{lemma}
\begin{proof}
For \(|z|>|q|\), define \(\mathcal F(z)=L((z-x)^{-1})\).
Its summand at \(u=q^m\) is
\[
 A\left(\frac1{z-u}-\frac1z\right)
 +\frac{(B+C)u}{(z-u)^2}+\frac{2Cu^2}{(z-u)^3}.
\]
The series is analytic at infinity. It converges locally normally
away from zero and the nodes, giving a meromorphic continuation to
\(\C\setminus\{0\}\). At each node the highest nonzero
weight supplies a genuine pole, of order three if \(C\ne0\), two
if \(C=0,B\ne0\), and one otherwise. Hence \(\mathcal F\) is nonrational.

If the size \(k\) determinant is zero, there is a nonzero polynomial
\(Q_k\) of degree less than \(k\) with \(L(x^jQ_k/D_k)=0\) for
\(0\le j<k\), where \(D_k(x)=\prod_{a=1}^k(1-q^a x)\).
We seek a polynomial \(T_k\) for which the error
\(\mathcal F Q_k-T_k\) vanishes at \(b,\ldots,b^k\) and is
\(O(1/z)\) at infinity. The factor \(D_k(z)\) gives the first
condition and the moment equations give the second. Put
\[
 E_k(z)= D_k(z)L_x\!\left(\frac{Q_k(x)}{(z-x)D_k(x)}\right),
 \qquad T_k(z)=\mathcal F(z)Q_k(z)-E_k(z).
\]
Combining fractions in the expression for \(T_k\) produces the
numerator \(Q_k(z)D_k(x)-D_k(z)Q_k(x)\), divisible by \(z-x\).
Thus \(T_k\) is a polynomial of degree less than \(k\).
The \(k\) moment equations cancel the first \(k\) terms of the expansion
of the interior transform at infinity, giving \(E_k(z)=O_k(z^{-1})\).
Indeed \(L\) is bounded in terms of the supremum of a function on
any fixed disk of radius between \(|q|\) and 1, by Cauchy's
inequalities and \eqref{eq:L}. The remainder after these \(k\)
terms is \(O_k(|z|^{-k-1})\) on that disk. The same boundedness
gives holomorphy of the interior transform at \(b,\ldots,b^k\).
Consequently \(E_k\) vanishes at those \(k\) points.

If two consecutive sizes are singular, the polynomial
\(T_{k+1}Q_k-T_kQ_{k+1}=E_kQ_{k+1}-E_{k+1}Q_k\) has degree at
most \(k-1\) by its growth at infinity, and vanishes at all \(k\)
common nodes. It is zero. If every sufficiently large size were
singular, all fractions \(T_k/Q_k\) in that range would therefore
represent one fixed reduced fraction \(T/Q\). Write
\(T_k=V_kT,\ Q_k=V_kQ\) with \(V_k\ne0\), allowing \(T=0,Q=1\).
The function \(E=\mathcal F Q-T\) is nonzero, meromorphic at infinity,
and \(V_k E=O_k(1/z)\). It follows that \(E(z)\sim c z^{-h}\) there
with \(h\ge1,\ c\ne0\), and \(\deg V_k\le h-1\).
Since \(|b^a|\to\infty\), the asymptotic gives \(E(b^a)\ne0\) for
all sufficiently large \(a\). At each such point with \(a\le k\),
\(E_k(b^a)=V_k(b^a)E(b^a)=0\), so \(V_k(b^a)=0\). The number of
these distinct zeros tends to infinity, contradicting the fixed
degree bound.
\end{proof}

For an integer \(b\ge2\), a rational dependence gives integers
\(A,B,C,\xi\) with \((A,B,C)\ne0\). Lemma~\ref{lem:clear} then
makes \(P_K(b,A,B,C,\xi)\) an integer. Its multiplier
\(b^{-e_K}H_K(b)\) is nonzero, so Lemma~\ref{lem:nonzero} gives a
nonzero integer for arbitrarily large \(K\). Lemma~\ref{lem:decay}
makes its absolute value less than one for all sufficiently large \(K\).
This contradiction and \eqref{eq:triangle} prove Theorem~\ref{thm:main}.

\section{A consequence for algebraic integers}

For an algebraic base \(b\), every partial sum of \(M_j(b)\) lies
in \(\Q(b)\). A failure of independence would give coefficients
\(A,B,C\) in this field, not all zero, for which the combination
\(\xi=AZ_1(b)+BZ_2(b)+CZ_3(b)\) also lies in \(\Q(b)\). This value
then has a conjugate at each embedding, even when the series at the
conjugate base diverge. We evaluate the polynomial \(P_K\) at the
conjugate parameters and bound their contribution to the field norm.

\begin{corollary}\label{cor:field}
Let \(b\) be an algebraic integer with \(|b|>1\). Suppose all of
its conjugates other than \(b\) and \(\bar b\) have modulus at most
one. Then \(1,M_1(b),M_2(b),M_3(b)\) are linearly independent
over \(\Q(b)\). This includes all real Pisot and Salem bases.
\end{corollary}
\begin{proof}
We first bound \(P_K(t,A,B,C,Y)\) on \(|t|\le1\) for any fixed
complex \(A,B,C,Y\). We give the details of the circle argument
from \cite[\S4]{Golden} because the conjugate of \(\xi\) at such a
place is simply a fixed number \(Y\).
On \(|t|=1+1/K\), the nodes \(r_a=t^a\) have modulus at most \(e\),
and \(|1-r_a/r_c|\ge1/(K+1)\) for \(a\ne c\). Thus
\[
 1/D(x)=\sum_{a=1}^K\frac{w_a}{1-x/r_a},\qquad
 |w_a|=\prod_{c\ne a}|1-r_a/r_c|^{-1}\le(K+1)^{K-1}.
\]
Each tail \(S_a\) is \(O_{A,B,C}(K^4)\) and \(\mu_j=O_{A,B,C}(K^3)\)
for \(0\le j\le2K\), since \(|t^m-1|\ge1/K\) for \(m\ge1\).
For \(U_{a,j}=\cl_{t,Y}(x^j/(1-x/r_a))\) the exact recurrence
\[
 U_{a,0}=Y-S_a,\qquad U_{a,j+1}=r_a(U_{a,j}-\mu_j)
\]
bounds \(|U_{a,j}|\) by \(\exp(O_{A,B,C,Y}(K+\log(2K)))\) for
\(j\le2K\). Summing \(w_a U_{a,j}\) and expanding the determinant
gives \(\log^+|\Delta_K|=O_{A,B,C,Y}(K^2\log(2K))\).

On that circle \(1/(K+1)\le |1-t^{-h}|\le2\) for \(h\ge1\).
M\"obius inversion bounds
\(\log|\Phi_d(t)|\le\ph(d)\log(1+1/K)+\tau(d)\log(2K+2)\),
where \(\tau(d)\) counts the divisors. In \eqref{eq:clearer},
\(c_d=O(K^2/d)\), and
\(\sum_{d\le K}\tau(d)/d\le(1+\log K)^2\).
Lemma~\ref{lem:clear} therefore gives
\(\log^+|H_K|=O(K^2\log^3(2K))\). Also \(|t^{-e_K}|\le1\).
The maximum modulus principle for the integral polynomial now yields
\[
 \max_{|t|\le1}\log^+|P_K(t,A,B,C,Y)|
                  =O_{A,B,C,Y}(K^2\log^3(2K))=o(K^3).
\]

Scale any putative field relation by one integer so that its
coefficients \(A,B,C\) and value \(\xi\) are algebraic integers.
For the unbounded sequence of nonzero sizes, \(P_K(b,A,B,C,\xi)\)
then has a nonzero integer norm. At the specified embedding,
Lemma~\ref{lem:decay} gives exponential decay on the \(K^3\) scale.
If \(b\) is nonreal, the complex conjugate of this embedding gives
the same absolute value. Each remaining embedding costs only
\(o(K^3)\) in the logarithm by the disk bound at its fixed conjugate
parameters. The absolute norm tends to zero on this sequence,
a contradiction. The change of variables \eqref{eq:triangle}
also holds by absolute convergence for complex \(|b|>1\).
\end{proof}

\section*{Verification}
Another AI model independently audited the mathematics and exposition
of the 18 November 2026 draft. It verified Theorem~\ref{thm:main}
and Corollary~\ref{cor:field}, rederived the decay margin, and
reproduced the finite checks with independently written code.

\small
\begin{thebibliography}{9}
\setlength{\itemsep}{1pt}
\bibitem{PVA} K. Postelmans and W. Van Assche, Irrationality of
\(\zeta_q(1)\) and \(\zeta_q(2)\), \emph{J. Number Theory}
\textbf{126} (2007), 119--154. \href{https://arxiv.org/abs/math/0604312}{arXiv:math/0604312}.
\bibitem{ZudPair} W. Zudilin, Approximations to \(q\)-logarithms
and \(q\)-dilogarithms, with applications to \(q\)-zeta values,
\emph{J. Math. Sci. (N.Y.)} \textbf{137} (2006), 4673--4683.
\href{https://www.math.ru.nl/~zudilin/PS/qhyp.pdf}{Theorem 2}.
\bibitem{Uplift} uplift-universal-inc, Irrationality of the Lambert-series
value \(\zeta_{1/p}(3)\), research draft, 12 October 2026.
\href{https://github.com/uplift-universal-inc/qzeta_3-irrationality/blob/main/paper/main.pdf}{Manuscript}.
\bibitem{UpliftCompanion} uplift-universal-inc, Results reserved for
companion articles, entry FUT-LI.
\href{https://github.com/uplift-universal-inc/qzeta_3-irrationality/blob/main/docs/FUTURE_RESULTS.md}{Research inventory}.
\bibitem{Kratt} C. Krattenthaler, A determinant identity for moments
of orthogonal polynomials that implies Uvarov's formula for the
orthogonal polynomials of rationally related densities.
\href{https://arxiv.org/abs/2103.03969}{arXiv:2103.03969}.
\bibitem{ZudHankel} W. Zudilin, On the irrationality of generalized
\(q\)-logarithm, \emph{Res. Number Theory} \textbf{2} (2016), article 15.
\href{https://doi.org/10.1007/s40993-016-0042-x}{doi:10.1007/s40993-016-0042-x}.
\bibitem{BZ} P. Bundschuh and W. Zudilin, Irrationality measures for
certain \(q\)-mathematical constants, \emph{Math. Scand.} \textbf{101}
(2007), 104--122. \href{https://doi.org/10.7146/math.scand.a-15034}{doi:10.7146/math.scand.a-15034}.
\bibitem{Golden} twinheterodox, The weight-three Lambert sum at the
golden ratio is irrational, Hexagon submission 2611.00623v1 (2026),
\S\S3--5.
\href{https://github.com/twinheterodox/algebraic-base-lambert-values/blob/main/output/pdf/algebraic-base-lambert-values.pdf}{Manuscript PDF}.
\bibitem{MersenneCode} twinheterodox, Verification programs for
the reciprocal Mersenne cubes (2026).
\href{https://github.com/twinheterodox/algebraic-base-lambert-values/tree/main/checks/mersenne}{Code and reference results}.
\end{thebibliography}
\end{document}
